% compile with:
% makeTex relazione2025
\documentclass[a4paper,12pt,twoside]{article}
\usepackage{latexsym}
\usepackage{amssymb}
\usepackage{amsmath}
\usepackage{siunitx}
\usepackage[version=4]{mhchem} % formule chimiche
\usepackage{cite}
\usepackage[resetlabels]{multibib}
\newcites{main}{10 pubblicazioni scelte dal 1 Novembre 2022 al 28 Febbraio 2025}
\newcites{all}{Lista di tutte le pubblicazioni dal 1 Novembre 2022 al 28 Febbraio 2025}
\usepackage[english,italian]{babel}
\pagestyle{headings}
\title{dfdfd}
%\maketitle
\usepackage[dvips]{graphicx}
\usepackage{eurosym}
\usepackage{fontenc}
\usepackage{epic,rotating,epsfig}
\usepackage{times}
\usepackage{fancyhdr}
\linespread{1.2}
%%%%%%%%%%%%%%%%%%%%%%%%%%%
\hoffset = -1.0 cm
\voffset = -2.0 cm
\textheight = 24.0 cm
\marginparwidth = 1.0 cm
\evensidemargin = 1.0 cm
\textwidth =16.1 cm
%%%%%%%%%%%%%%%%%%%%%%%%%%%
\begin{document}
\fontencoding{OT1}\fontfamily{pbk}\fontseries{m}
\selectfont
\vspace*{2.20cm}

\begin{center}
  {\LARGE \bf Relazione sull'attivit\`a svolta dal 1 Novembre 2022 al 28 Febbraio 2025} \\
  \vspace*{0.2cm}
  {\large \it Emanuele Di Marco}
\end{center}

\vspace*{2.2cm}

\hspace*{-0.6cm}{ L'attivit\`a da me svolta in questo periodo si inserisce nella ricerca di fenomeni non
  previsti dal Modello Standard delle particelle elementari
  nell'ambito dei dei fenomeni elettrodeboli e del bosone di Higgs,
  nonch\'e di misure di parametri fondamentali dello stesso,
  soprattutto nell'ambito dell'esperimento {\it CMS} al {\it Large
    Hadron Collider} (LHC) del {\it CERN} di Ginevra. Allo stesso
  tempo ho contribuito alle misure relative all'esperimento {\it
    CYGNO}, per la ricerca di segnali di materia oscura, utilizzando
  tutti i dati presi dal primo prototipo, {\it LIME}, il quale \`e
  stato installato, commissionato e ha preso dati in maniera
  continuativa nei {\it Laboratori Nazionali del Gran Sasso} (LNGS).

  In questo periodo, nell'ambito dell'attivit\`a dell'esperimento CMS,
  mi sono dedicato a studi sul bosone di Higgs. Dalla primavera del
  2022, infatti, l'acceleratore LHC ha fornito collisioni
  protone-protone ($pp$) ad un centro di massa maggiore di quello
  della presa dati del Run-2: $\sqrt{s}=\SI{13.6}{TeV}$, e
  l'esperimento CMS ha raccolto circa \SI{41.5}{fb^{-1}},
  \SI{32.7}{fb^{-1}} e \SI{122.2}{fb^{-1}} nel 2022, 2023 e 2024,
  rispettivamente. Quindi da un lato \`e stato possibile misurare le
  sezioni d'urto di produzione del bosone di Higgs a questa energia
  del centro di massa mai raggiunta prima, e confrontarle con le
  predizioni teoriche. Dall'altro lato, con questa luminosit\`a
  integrata, confrontabile a quella raccolta con i precedenti run di
  LHC, \`e stato possibile fare nuove misure rese accessibili dal
  campione di dati due volte maggiore, e dalla maggiore sezione d'urto di
  produzione a \SI{13.6}{TeV}. Sono uno degli autori principali, nonch\'e 
  supervisor di due studenti di dottorato di cui sono relatore, di due analisi attualmente in fase di
  pubblicazione. 
  
  La prima \`e la misura della larghezza intrinseca del bosone di
  Higgs, che essendo prevista dal Modello Standard essere
  $\Gamma_H=\SI{4.1}{MeV}$, \`e di circa tre ordini di grandezza
  inferiore rispetto alla risoluzione sperimentale dei canali con alta risoluzione di
  massa, $H\to ZZ\to 4\ell$ e $H\to\gamma\gamma$. Questa
  \`e una misura innovativa che sfrutta lo shift di massa indotto
  dall'interferenza quantistica tra l'ampiezza risonante e quella del
  processo non risonante $pp\to\gamma\gamma$. La sensibilit\`a \`e
  quella di una decina di MeV, un ordine di grandezza maggiore al limite superiore
  ottenuto dagli eventi {\it off-shell} di $H\to
  4\ell$, ma con assunzioni teoriche molto minori e comunque un ordine di grandezza migliore del vincolo dalla misura diretta.  La seconda analisi dei dati
  di CMS di cui sono autore principale insieme a una studentessa di dottorato di cui sono relatore, sfrutta anch'essa il decadimento $H\to\gamma\gamma$, e
  ricerca possibili accoppiamenti anomali del bosone
  di Higgs ai bosoni vettori ($V=W,Z$), tramite la produzione VBF e
  associata $VH$.

 Dal 2021 fino a Settembre 2023 sono stato il chair {\it Publication
   Committee} del gruppo del bosone di Higgs di CMS. Nel corso di
 questo periodo ho coordinato le pubblicazioni di tutti gli articoli
 (circa 30) sul bosone di Higgs. Dal 1 Settembre 2023 ad oggi sono il
 convener dell'intero gruppo di fisica dell'Higgs, coordinando il
 lavoro di circa 300 persone coinvolte nelle relative analisi dei
 dati. Durante questo periodo ho condotto ai risultati preliminari, o
 alla pubblicazione su rivista, circa 50 analisi. La mia
 attivit\`a in questo ambito va dalla pianificazione dei risultati di
 fisica delle linee di ricerca sull'Higgs sia nell'immediato futuro, che con il dataset completo di LHC prima dell'upgrade alla fase 2 (Run2+Run3, previsto terminare la presa dati
 a met\`a del 2026). Essa include la supervisione dei sottogruppi ($H\to\gamma\gamma$,
 $H\to ZZ$, $H\to WW$, $H\to b\bar{b}$, $H${\it LepRare}), nel quale l'aspetto fondamentale \`e quello della guida e revisione di tutte le analisi, dall'impostazione iniziale della strategia, alla comprensione dei risultati, fino alla fase di pubblicazione su rivista. Una parte considerevole dell'impegno \`e quello del coordinamento tra il gruppo di analisi e quelli di {\it offline
   and computing} e delle performance dei rivelatori strettamente legate
 ai risultati di fisica. 
 
 I risultati in cui ho dato un contributo personale determinante, e pubblicate nel periodo considerato, sono: la combinazione ATLAS+CMS che ha
 permesso la prima evidenza del decadimento raro $H\to
 Z\gamma$~\citemain{hzg}, i vincoli sull'auto-accoppiamento del bosone
 di Higgs dalla combinazione delle misure di singolo $H$ e doppio
 $H$~\citemain{HplusHH}, la prima osservazione del decadimento $H\to b\bar{b}$
 con la produzione $VH$~\citemain{VHbb}, la prima misura della produzione
 rara $bbH$, attraverso i decadimenti $H\to\tau\tau$ o $H\to WW$
 ~\citemain{bbH}, la prima misura della produzione $ttH(b\bar{b})$,
 includendo anche stati finali completamente adronici (sottomessa a rivista, in fase di pubblicazione). Lo sviluppo di
 algoritmi di {\it tagging} basati su Machine Learning, come {\it
   Graph Neural Networks}, atti a determinare il sapore dei quark
 originanti i jet adronici, ha permesso di determinare il limite superiore pi\`u basso
 per l'accoppiamento di Yukawa del bosone di Higgs ai quark di seconda
 generazione, $H\to c\bar c$, attraverso la produzione associata
 $VH$~\citemain{VHcc}. L'applicazione di queste tecniche, applicate
 agli stati finali con alto boost di Lorentz del bosone di Higgs, ha
 reso accessibili alcuni decadimenti completamente adronici, che fino
 ad oggi erano al di l\`a della portata sperimentale, permettendo di
 abbattere di ordini di grandezza il fondo di multijet da QCD. Esempi
 su cui il mio contributo nella definizione della strategia di analisi
 e revisione in qualit\`a di convener dell'Higgs ha permesso la
 pubblicazione in questo periodo sono: $H\to \tau\tau$~\citemain{boostedHtautau}
 e la produzione in gluon-fusion e VBF di $H\to b\bar{b}$~\citemain{boostedHbb}, in un regime
 in cui i prodotti di decadimento adronici (dei $\tau$ o dei quark
 $b$) sono inglobati in uno stesso jet ricostruito a causa dell'alto
 boost. L'analisi della sotto-struttura e delle correlazioni dei
 prodotti dell'adronizzazione sono risultati fondamentali, ma ha
 richiesto un lavoro raffinato nella modellizzazione
 dell'adronizzazione con la simulazione, e inoltre calibrazioni
 precise della risposta di questi algoritmi tra dati e
 simulazioni. Queste tecniche rendono possibile ampliare la sensibilit\`a a meccanismi rari di produzione del bosone di Higgs, e in particolare quello del doppio-$H$ non risonante, che offre il metodo migliore per vincolare l'accoppiamento $H-H$, il quale determina la forma del potenziale del campo di Higgs, che \`e correlato alla stabilit\`a quantistica dell'Universo. Insieme a un terzo dottorando, di cui sono relatore, mi sto cominciando ad occupare di $HH\to b\bar{b}\tau\tau$, uno dei 3 canali che, grazie alla combinazione di alto {\it Branching Ratio} e relativamente alto rapporto segnale/fondo, sar\`a fondamentale per determinare i limiti pi\`u stringenti con i dati della fase-1 di LHC. Personalmente sto iniziando l'analisi dei dati di Run3 di LHC, ma ho coordinato come convener del gruppo di Higgs l'analisi dei dati di Run2~\citemain{HHbbtautau}. 
 
 Inoltre, sebbene non ancora pubblicate, ma con risultati
 preliminari gi\`a presentati a conferenze internazionali, ho
 partecipato e coordinato le prime misure della sezione d'urto,
 fiduciale e differenziale, del bosone di Higgs a
 $\sqrt{s}=\SI{13.6}{TeV}$, nei canali $H\to\gamma\gamma$ e
 $H\to ZZ\to 4\ell$, nonch\'e molteplici altre misure di precisione che
 hanno permesso di vincolare effetti di nuova fisica a scale di
 energia inaccessibili direttamente a LHC attraverso le
 interpretazioni nell'ambito delle {\it Teorie di Campo Effettive}.

In questo periodo \`e stata completata, resa pubblica e sottomessa su
rivista ``Nature'', la misura della massa del bosone $W$,
$m_W=\SI{80360.2\pm 9.9}{MeV}$, in accordo con la predizione del
Modello Standard, e con la maggior precisione sperimentale mai
raggiunta. Questo risultato completa l'analisi dei dati di CMS del 2020
che era stata usata per la misura delle sezioni d'urto
differenziali, in stati di elicit\`a e di carica, su cui ho lavorato
come autore principale, negli anni precedenti al periodo in
considerazione, con la stessa tecnica di estrazione del segnale usata
per la misura di $m_W$. Essa infatti ha permesso di vincolare le
sistematiche delle PDF a LHC che erano uno dei fattori limitanti per $m_W$ a un
collisore $pp$. L'aggiunta della calibrazione di scala di momento dei muoni con una precisione di
$\approx 10^{-4}$ ha permesso quindi finalmente l'estrazione di $m_W$ con la precisione
di \SI{10}{MeV}. Esso \`e il primo passo verso una misura che
raggiunga la precisione teorica del fit elettrodebole di circa
\SI{5}{MeV}, che quindi potrebbe permettere un test della consistenza
del Modello Standard, e in particolare delle correzioni quantistiche
della massa del $W$: per raggiungerlo servir\`a analizzare tutti i
dati di Run2 con la stessa precisione, e riducendo ancora di pi\`u le
assunzioni teoriche sul modello di produzione del bosone $W$ a LHC.
 
Durante questo periodo, ho continuato ad essere
co-responsabile del centro di calcolo (\textit{Tier-2}) per i dati di
CMS di Roma, un sistema condiviso con l'esperimento ATLAS. Esso \`e
costituito da dieci rack che contengono qualche centinaio di server di
elaborazione dati (per un totale di circa 2500 core) e archiviazione
dei dati (per un totale di circa 2\,PB). Il centro di calcolo permette
l'eleborazione di circa il 20\% dei dati raccolti dal LHC ogni anno,
ed \`e connesso al sistema di \textit{grid} mondiale. In questi anni
mi sono occupato della pianificazione degli upgrade e degli acquisti
per il Tier-2 di CMS, per un investimento medio di circa 100\,k\euro{}
all'anno, nonch\'e dell'installazione delle risorse di calcolo e di
storage dati acquistate con i fondi relativi al {\it PNRR}.

Nel periodo di interesse ho continuato la partecipazione allo sviluppo
dell'esperimento CYGNO, un rivelatore per la ricerca diretta di
materia oscura di tipo \textit{Weak Interacting Massive Particle}
(WIMP), che si manifesti nei rinculi delle ipotetiche particelle sui
nuclei di atomi di gas, con energia cinetica di pochi keV. Nel periodo
2022-2025 il prototipo {\it LIME} \`e stato portato ai LNGS, 
e ho quindi partecipato sia al commissioning
del rivelatore che alle calibrazioni in energia, con il primo dataset
completo raccolto {\it underground}.  Fino a Settembre 2023 sono stato
il coordinatore del gruppo di analisi dati dell'esperimento CYGNO, e
mi sono occupato del codice e dell'infrastruttura per la ricostruzione
{\it prompt} dei dataset raccolti dall'esperimento, usata da tutti i
membri della collaborazione internazionale. Ho sviluppato
ulteriormente tali algoritmi, adattandoli all'{\it occupancy} media
pi\`u ridotta rispetto ai test con il prototipo in superficie,
soprattutto nella configurazione con uno {\it shield} di \SI{5}{cm} di
rame radiopuro intorno alla TPC. Sono l'autore principale delle misure
descritte nel'articolo~\citemain{LimePerf}, le quali hanno permesso di individuare aspetti specifici, quali quello dell'elettroluminescenza, descritto nell'articolo~\citemain{LimeElectro}.
 
Durante il periodo in considerazione mi sono occupato anche di
didattica relativa al mio campo di interesse, come esercitatore per il
corso di Fisica Nucleare e Subnucleare 1 nell'a.a. 2022-2023, come
tutor di esperienze di laboratorio per il corso di Laboratorio di
Fisica Nucleare 2 (dal 2019 ad oggi), come relatore di 3 tesi
magistrali, 3 tesi di dottorato ancora in corso e numerose
dissertazioni triennali.

Inoltre in questo periodo sono stato referee articoli scientifici per
la rivista \textit{PLB} (``Physics Letters B'', nel settore di ``Particle
Physics, Nuclear Physics and Cosmology''), per la rivista
\textit{JHEP} (``Journal of High Energy Physics''), per la rivista \textit{ROPP} (``Reports on Progress in Physics''), per la rivista \textit{MLST} (``Machine Learning: Science and
Technology''), ed editor della sezione ``Upgrades in High Energy
Physics Experiments'' della rivista \textit{Symmetry} di MDPI.

  
}

\bibliographystylemain{unsrt}
\bibliographymain{mainpub}

\citeall{*}
\bibliographystyleall{unsrt}
% trovato su INSPIRE con la query: "a emanuele di marco and de > 2023-11-01 and a"
\bibliographyall{INSPIRE-CiteAll-Nov2022-March2025}


\end{document}




